\section{Round-twelve positive closure: polymer characteristic pressure and sharp regular shells}
\label{sec:r12-b1}

The proof uses the full complex pressure of the dynamically tilted gas.  It
does not partition the fixed torus into linearly many causally separated
blocks.  Compact and high frequencies are controlled directly by the
singleton term and the absolutely convergent real-trajectory polymer
remainder.

\subsection{Pressure decomposition and strict complex aperiodicity}

For a path/contact source $H$ and constraint multiplier $\lambda$ write
\[
 Q_\varepsilon(H,\lambda)
 ={1\over\mu_\varepsilon}
 \log\mathbb E\exp\{\mu_\varepsilon H+
                    \lambda\cdot\textstyle\sum_i C(z_i(0))\}.
\]
The B2 grand-canonical expansion gives
\[
 Q_\varepsilon=Q_\varepsilon^{(1)}+R_\varepsilon,
 \qquad
 \sup_{\mathcal U}\|D^jR_\varepsilon\|\le C_jT_*,\quad j\le4,
\]
on a complex source--multiplier polydisc.  Here $Q^{(1)}$ is the exact
compound-Poisson singleton pressure and $T_*$ is chosen below.

\begin{lemma}[Uniform constraint covariance]
\label{lem:r12-b1-covariance}
Assume that the mark map $C$ has full affine span on one finite family of
positive-activity phase boxes.  There are $T_*>0$, $c>0$, and a common
polydisc such that
\[
 D^2_{\lambda\lambda}Q_\varepsilon(H,\lambda)\ge cI
\]
for every real $(H,\lambda)$ in the polydisc and all small $\varepsilon$.
\end{lemma}

\begin{proof}
The singleton Hessian is the activity integral of
$(C-\bar C)(C-\bar C)^T$ and is uniformly positive by the affine-span
hypothesis and positivity on the chosen boxes.  The connected Hessian is
bounded by $CT_*$ by the B2 trajectory-cluster norm.  Choose $T_*$ so that
this perturbation is at most half the singleton lower eigenvalue.
\end{proof}

Let $(t,u)\in[-\pi,\pi]\times\mathbb R^{d-1}$ be the number-lattice and
continuous frequencies.  At the exact finite saddle set
\[
 \Phi_\varepsilon(t,u)=
 Q_\varepsilon(H,\lambda_{H,a,\varepsilon}+i(t,u))
 -Q_\varepsilon(H,\lambda_{H,a,\varepsilon}).
\]

\begin{theorem}[Full-frequency pressure gap]
\label{thm:r12-b1-frequency}
There are $\delta,R,c>0$ and, for every $s>d+3$, a $C_s$ such that
\[
 \Re\Phi_\varepsilon(t,u)\le
 \begin{cases}
  -c(|t|^2+|u|^2),& |t|+|u|\le\delta,\\
  -c,& \delta\le |t|+|u|\le R,\\
  -c-C_s\log(1+|u|),& |u|>R,
 \end{cases}
\]
uniformly on compact source and target charts.  Equivalently the full
characteristic function is Gaussian near zero, exponentially small in
$\mu_\varepsilon$ on every compact nonzero annulus, and bounded by
\[
 e^{-c\mu_\varepsilon}(1+|u|)^{-s\mu_\varepsilon}
\]
in the continuous tail.  The three regions cover the complete Fourier
fundamental domain.
\end{theorem}

\begin{proof}
Near zero use Lemma~\ref{lem:r12-b1-covariance} and Taylor expansion.  On a
compact annulus, equality in the triangle inequality for the tilted singleton
integral would make $(t,u)\cdot C$ constant modulo $2\pi$ on every full-rank
phase box.  The number span and continuous affine span force the trivial
frequency.  Compactness gives a negative singleton gap; the $O(T_*)$ polymer
remainder is absorbed by the choice of $T_*$.  For large $u$, integrate the
smooth singleton activity on the full-rank boxes $s$ times.  Its Fourier
transform is $O((1+|u|)^{-s})$, while the Poisson exponent contains
$\mu_\varepsilon$ independent singleton opportunities at the logarithmic
pressure level.  Normal convergence of connected polymers changes the real
part by at most $CT_*$ and does not undo the negative activity mass.  This
gives the final speed-dependent power.
\end{proof}

\subsection{The finite saddle and its mean image}

\begin{theorem}[Uniform finite-volume mean map]
\label{thm:r12-b1-meanmap}
There is a bounded convex multiplier domain $\Lambda$ and a compact target
set $A$ in the relative interior of the constraint convex support such that,
for every $H$ in the source chart and small $\varepsilon$,
\[
 \lambda\mapsto D_\lambda Q_\varepsilon(H,\lambda)
\]
is a degree-one diffeomorphism from $\Lambda$ onto a neighborhood of $A$.
The exact saddle $\lambda_{H,a,\varepsilon}$ and all derivatives up to order
three converge locally uniformly to their limiting values.
\end{theorem}

\begin{proof}
Strict convexity gives local injectivity.  Choose $\Lambda$ as a smooth
sublevel of the singleton Legendre function containing the limiting saddles
for $A$.  On its boundary the singleton gradient points strictly outward
relative to $A$; the connected pressure is a $CT_*$ perturbation, so the same
holds for the full finite pressure.  Brouwer degree is one and no two
preimages can occur by strict convexity.  Normal complex convergence and the
implicit-function theorem give derivative convergence.
\end{proof}

\subsection{Regular shell coefficients}

Let $W_\varepsilon\subset\mathbb R^{d-1}$ be centered Lipschitz domains in
CLT coordinates satisfying
\[
 \gamma_\Sigma(W_\varepsilon)\ge\mu_\varepsilon^{-M},
 \qquad
 \gamma_\Sigma((\partial W_\varepsilon)^{+h})
 \le C h\mu_\varepsilon^M
\]
for $0<h<1$, and let the number coordinate be exact.  Centered expanding boxes
with $\sqrt\mu\delta_\varepsilon\to\infty$ are included.

\begin{theorem}[Relative mixed local coefficient]
\label{thm:r12-b1-coefficient}
Uniformly on the source and target chart,
\[
 \begin{aligned}
 &\mathbb P_{H,\lambda_{H,a,\varepsilon}}
 \left(N=N_\varepsilon,
 \sqrt{\mu_\varepsilon}(Y-a)\in W_\varepsilon\right)\\
 &\qquad={\gamma_{\Sigma^c_{H,a}}(W_\varepsilon)\over
 \sqrt{2\pi\mu_\varepsilon\sigma^2_{N\mid Y}}}
 \left(1+o(1)\right),
 \end{aligned}
\]
where the error is relative to the displayed Gaussian mass.  In particular
its logarithm is $o(\mu_\varepsilon)$.
\end{theorem}

\begin{proof}
Fourier inversion is centered at the exact saddle.  On a ball of radius
$\mu^{-2/5}$, the fourth-order pressure bound gives a relative Edgeworth
error.  Smooth upper and lower approximations of $1_{W_\varepsilon}$ use a
boundary thickness $h=\mu^{-K}$; the declared Minkowski estimate makes their
Gaussian difference $o(\gamma(W_\varepsilon))$.  The compact annulus and tail
are $e^{-c\mu}$ times an integrable speed-dependent power by
Theorem~\ref{thm:r12-b1-frequency}, hence are
$o(\mu^{-1/2}\gamma(W_\varepsilon))$.  Division by the Gaussian main term
gives the relative formula.
\end{proof}

\subsection{Source-dependent microcanonical transfer}

\begin{theorem}[Prepared pressure and initial covariance]
\label{thm:r12-b1-main}
For the primitive shell in Theorem~\ref{thm:r12-b1-coefficient},
\[
 Q^{\rm mc,a}(H)=
 Q(H,\lambda_{H,a})-\lambda_{H,a}\cdot a
 -Q(0,\lambda_{0,a})+\lambda_{0,a}\cdot a.
\]
The initial empirical law has the constrained entropy rate, and its Gaussian
covariance is
\[
 \Sigma^{\rm mc}=\Sigma-
 \Sigma C^*(C\Sigma C^*)^{-1}C\Sigma.
\]
The theorem applies jointly to the B2 particle-path and actual-contact source
class.
\end{theorem}

\begin{proof}
Change measure in numerator and denominator at their separate exact finite
saddles.  Theorem~\ref{thm:r12-b1-coefficient} shows that both residual shell
probabilities have subexponential logarithm and supplies their relative
coefficients.  Passing to the normal pressure limit gives the displayed
formula.  Convex duality gives the initial rate.  Differentiating the saddle
equation twice gives the Schur complement covariance.  The only dynamical
input is the earlier B2 grand-canonical pressure theorem, so the dependency is
noncircular.
\end{proof}
